%%=============================================================================
%% DADOS DO TRABALHO -- FONTE ÚNICA DA VERDADE
%%
%% Este é o ÚNICO arquivo onde os dados do seu trabalho devem ser digitados.
%% Tudo o que aparece na capa, folha de rosto, ficha catalográfica, folha de
%% aprovação e nos metadados do PDF é gerado a partir daqui.
%%
%% Campos marcados como OBRIGATÓRIO que ficarem em branco aparecem no PDF
%% como "[[ PREENCHER: campo ]]" e geram aviso na compilação. Ao gerar a
%% versão final, use a opção `entrega' em main.tex: ela transforma o aviso
%% em erro e impede a compilação com campos pendentes.
%%
%% Para reaproveitar estes dados no corpo do texto, use os acessores:
%%   \TituloTrabalho  \NomeAutor      \NomeOrientador   \NomeCoorientador
%%   \NomeCurso       \NomePrograma   \NomeInstituicao  \AreaConcentracao
%%   \LinhaPesquisa   \LocalTrabalho  \CidadeTrabalho   \AnoTrabalho
%% Exemplo: "Ao meu orientador, \NomeOrientador, agradeço..."
%%=============================================================================

%% --- Identificação ---------------------------------------------------------
\titulo{Inteligência artificial na triagem de pendências repetitivas na Advocacia Pública}  % OBRIGATÓRIO
\subtitulo{o caso da PGDF}                        % opcional: \subtitulo{}
\titulotraduzido{Artificial intelligence in the triage of repetitive pending tasks in public legal advocacy: the case of PGDF}

\autor{Analécia Hanel Rorato}                     % OBRIGATÓRIO
\autorficha{RORATO, Analécia Hanel}               % OBRIGATÓRIO (formato invertido)

\orientador{Prof. Dr. Marcelo Rodrigo de Souza Pita}  % OBRIGATÓRIO
\coorientador{}                                     % opcional

%% --- Vínculo institucional -------------------------------------------------
\curso{}                                 % OBRIGATÓRIO em bacharelado/especialização
\programa{Administração Pública}         % OBRIGATÓRIO em mestrado/doutorado
\area{Administração Pública}             % área de concentração (pós-graduação)
\linhapesquisa{Ciência de Dados e Inteligência Artificial}  % linha de pesquisa

\cidade{Brasília}                        % usada na folha de aprovação
\local{Brasília-DF}                      % usado na capa e folha de rosto
\data{2027}                              % ano de depósito

%% --- Defesa ----------------------------------------------------------------
%% Informe APENAS a data; a cidade vem de \cidade acima.
\datadefesa{7 de maio de 2027}           % OBRIGATÓRIO

%% Banca: \membrobanca[Instituição]{Nome}{Função}   (instituição é opcional)
\membrobanca{Prof. Dr. Marcelo Rodrigo de Souza Pita}{Orientador}
\membrobanca{Prof. Dr. Hélio Bomfim de Macêdo Filho}{Examinador}
\membrobanca{A definir}{Examinador}

%% --- Ficha catalográfica ---------------------------------------------------
%% A ficha OFICIAL é elaborada pela Biblioteca Ministro Moreira Alves.
%% Solicite por biblioteca@idp.edu.br e transcreva os dados recebidos aqui.
%% Os valores abaixo são provisórios, até a Biblioteca enviar a ficha.
\cutter{}                                % notação Cutter fornecida pela Biblioteca
\numeropaginas{18}                       % OBRIGATÓRIO: total de folhas
\cdd{351}                                % OBRIGATÓRIO: provisório (351 = Administração pública)
\assuntos{1. Inteligência artificial. 2. Advocacia Pública. 3. Administração Pública. I. Título.}

%% --- Resumo ----------------------------------------------------------------
\palavraschave{Inteligência artificial. Advocacia Pública. Triagem processual. Governança algorítmica. Supervisão humana.}   % OBRIGATÓRIO
\keywords{Artificial intelligence. Public legal advocacy. Case triage. Algorithmic governance. Human oversight.}             % OBRIGATÓRIO

%% --- Sobreposições opcionais -----------------------------------------------
%% Texto da natureza do trabalho, na redação do programa.
\preambulo{Dissertação apresentada como requisito parcial para obtenção do
           título de Mestre em Administração Pública do Mestrado Profissional
           em Administração Pública, Ciência de Dados e Inteligência
           Artificial, do Instituto Brasileiro de Ensino, Desenvolvimento e
           Pesquisa (IDP).}

%% Tipo do trabalho na ficha catalográfica (padrão: Monografia/Dissertação/Tese)
% \tipotrabalho{Trabalho de Conclusão de Curso}

%% Linha inteira da data de defesa, se o seu programa exigir outra redação
% \datadefesacompleta{Aprovada em 15 de dezembro de 2026.}

%% Logotipo da capa
% \logo{figuras/idp-logo.png}
% \semlogo
